% !TEX TS-program = pdflatex
% CV réseau LinkedIn — Entreprises Cybersécurité pure — SOC/Pentest/WAF
\documentclass[a4paper,9pt]{article}
\usepackage[utf8]{inputenc}
\usepackage[T1]{fontenc}
\usepackage[scaled=0.90]{helvet}
\renewcommand{\familydefault}{\sfdefault}
\usepackage[left=0.8cm,right=0.8cm,top=0.3cm,bottom=0.25cm]{geometry}
\usepackage{enumitem}
\usepackage{titlesec}
\usepackage{xcolor}
\usepackage[hidelinks]{hyperref}
\definecolor{navy}{RGB}{23,54,93}
\pagestyle{empty}
\setlength{\parindent}{0pt}
\setlength{\parskip}{0pt}
\linespread{0.82}
\hypersetup{pdftitle={CV - Baha Eddine Belhaj Mouldi - Cybersecurite SOC Pentest WAF},pdfauthor={Baha Eddine Belhaj Mouldi},pdfsubject={Cybersecurite, SOC, Pentest, WAF, Machine Learning}}
\titleformat{\section}{\normalsize\bfseries\color{navy}}{}{0pt}{\MakeUppercase}[{\vspace{1pt}\color{navy}\hrule height 0.6pt}]
\titlespacing*{\section}{0pt}{1.5pt}{0.8pt}
\setlist[itemize]{leftmargin=1.0em,label=\textbullet,itemsep=0pt,topsep=0pt,parsep=0pt}
\newcommand{\cventry}[4]{\textbf{#1} \hfill \textbf{#4}\\[0.2pt]#2{\ifx\relax#3\relax\else, #3\fi}\par\vspace{0.3pt}}
\begin{document}
\begin{center}
{\LARGE\bfseries\color{navy} Baha Eddine Belhaj Mouldi}\\[2pt]
{\normalsize\color{navy} Ingénieur Cybersécurité \textbar\ WAF \& ML \textbar\ Tests d'intrusion \textbar\ Détection de menaces}\\[3pt]
{\small Tunis, Tunisie \quad|\quad +216 55 128 982 \quad|\quad \href{mailto:bahaeddine.belhajmouldi@gmail.com}{bahaeddine.belhajmouldi@gmail.com}}\\[1pt]
{\small \href{https://www.linkedin.com/in/bahamouldi}{linkedin.com/in/bahamouldi} \quad|\quad \href{https://github.com/bahamouldi}{github.com/bahamouldi} \quad|\quad \href{https://bahamouldi.github.io/portfolio/}{bahamouldi.github.io/portfolio}}
\end{center}
\vspace{1pt}
\section{Profil}
Ingénieur cybersécurité (ENIT, mention Très Bien), auteur de BeeWAF : un pare-feu applicatif intelligent combinant moteur ML (Random Forest, Gradient Boosting), 10 041 règles de détection et détection d'évasion à 18 couches, testé sur 153 scénarios OWASP avec Burp Suite Professional. Expérience de cycle SOC complet (détection, investigation, réponse à incident) et de gestion des vulnérabilités (virtual patching de dizaines de CVE). Passionné par la recherche offensive autant que défensive --- de la conception de règles de détection à leur contournement pour les renforcer. Rigueur, autonomie et goût pour les défis techniques.
\section{Compétences clés}
\textbf{Sécurité applicative \& WAF} : conception de moteurs de détection (règles + ML), OWASP Top 10/API/LLM, détection d'évasion, virtual patching de CVE\\[0.3pt]
\textbf{Tests d'intrusion} : Burp Suite Professional, CVSS/CWE, pentest web/API, 153 scénarios d'attaque testés\\[0.3pt]
\textbf{SOC} : SIEM (Elastic Stack, Kibana), règles Sigma, TheHive, réponse à incident (détection $<$2 min, blocage $<$7 min)\\[0.3pt]
\textbf{Machine Learning appliqué à la sécurité} : Random Forest, Gradient Boosting, Isolation Forest, évaluation de modèles (CSIC 2010)\\[0.3pt]
\textbf{Systèmes \& Réseaux} : Windows Server, Active Directory, Linux (Kali, Ubuntu, CentOS, Debian), VPN, Proxy, Firewall\\[0.3pt]
\textbf{Outillage} : Python, Bash, PowerShell, Docker, Kubernetes, CI/CD, N8N
\section{Expérience professionnelle}
\cventry{Ingénieur Sécurité --- Stagiaire PFE (mention Très Bien)}{Beehive Entreprises}{Tunisie}{02/2026 -- 06/2026}
\begin{itemize}
  \item Conception de BeeWAF : pipeline à 27 modules, 3 modèles ML, détection d'évasion 18 couches, testé sur 153 scénarios OWASP (Burp Suite Professional).
  \item Score 98,2/100, 0\% faux positifs ; virtual patching de 37 CVE ; conformité continue sur 7 référentiels.
\end{itemize}
\cventry{Spécialiste Cybersécurité, Freelance (1 an)}{Beehive Entreprises}{Tunisie}{01/2025 -- 12/2025}
\begin{itemize}
  \item Tests d'intrusion et audits de vulnérabilités sur 8 applications web/API, rapports de remédiation.
\end{itemize}
\cventry{Stagiaire Cybersécurité --- Détection \& Réponse à incident SAP}{Streamlink}{Tunisie}{07/2025 -- 08/2025}
\begin{itemize}
  \item Pipeline SOC automatisé (Python/PyRFC, ELK, N8N, TheHive) : détection $<$2 min, blocage $<$7 min.
\end{itemize}
\cventry{Stagiaire Sécurité Réseau}{Tunisie Telecom}{Tunisie}{07/2024}
\begin{itemize}
  \item Détection SOC (Elastic/Kibana, règles Sigma) : 10+ vulnérabilités identifiées.
\end{itemize}
\section{Projets clés}
\cventry{BeeWAF --- WAF Intelligent (PFE, 86/100)}{FastAPI, Python, ML, Docker, Kubernetes, ELK}{}{2026}
\begin{itemize}
  \item Pipeline à 27 modules, 107 signatures, détection d'évasion 18 couches. Latence $<$20ms, 100\% attaques OWASP Top 10 bloquées en test.
\end{itemize}
\cventry{Système de Détection et Réponse aux Incidents SAP}{Python, PyRFC, ELK, N8N, TheHive}{}{2025}
\begin{itemize}
  \item Pipeline SOC complet : collecte $\to$ corrélation $\to$ réponse automatisée sur 3 scénarios d'incident.
\end{itemize}
\section{Formation}
\cventry{Diplôme d'Ingénieur en Génie Informatique --- Mention Très Bien}{École Nationale d'Ingénieurs de Tunis (ENIT)}{Tunisie}{09/2023 -- 06/2026}
\begin{itemize}
  \item Diplômé le 25/06/2026. Classé 65e sur 600+ au concours national d'entrée.
\end{itemize}
\cventry{Classes préparatoires Physique-Technologie}{Institut Préparatoire aux Études d'Ingénieurs, El Manar}{Tunisie}{09/2021 -- 06/2023}
\section{Certifications}
Réseaux (CCNA) --- Cisco (2025) ; Fondements du Deep Learning --- NVIDIA (2025) ; IA Adversariale --- NVIDIA (2025) ; Machine Learning --- NVIDIA (2024) ; Git \& GitHub --- 365 Data Science (2024).
\section{Langues}
Français (Courant) \quad|\quad Anglais (B2) \quad|\quad Arabe (Natif)
\end{document}
